% Objectives
\section{Objectives}

%We will use requests per second as our measure of throughput for the overall system as it can be used as an overall measure of performance
%
%encouraging a cost effective solution


\subsection{Throughput Modelling}

Demonstrated in \autoref{chap:simulator-usage}, the Shallowsim simulator has support for estimating the throughput of a single request of a static input sequence and decode length.
It is highly improbable that an entire workload consists of a single request type, instead we expect a distribution of sequence lengths, as shown in \autoref{chap:workload-modeling}.
Hence, we need a method to accurately calculate the throughput of a system, with a variety of input sequence.

Using the probability distribution of the workload, we can compute the weighted average of throughput, using measured performance values for each sequence length.
Let $R_{\text{prefill}}(s)$ denote the measured prefill throughput (in requests per second) for a request of sequence length $s$.
The expected prefill throughput for workload $w$, denoted $\mathbb{E}[R_{\text{prefill}} \mid w]$ (expectation), is then given by:

\[
\mathbb{E}[R_{\text{prefill}} \mid w] = \sum_{s \in S} R_{\text{prefill}}(s) \cdot p(s \mid w)
\]

In addition to prefill, we can also compute the expected throughput for the decode phase. 
While prefill throughput depends solely on the input sequence length $s$, decode throughput typically depends on both the input length $s$ and the output (decode) length $d$.
For simplicity, we assume a fixed average decode length $d$ as a parameter.
Let $R_{\text{decode}}(s, d)$ denote the measured decode throughput (in requests per second) for an input sequence length $s$ and average decode length $d$.
Given a workload-specific sequence length distribution $p(s \mid w)$, the expected decode throughput for workload $w$ is:

\[
\mathbb{E}[R_{\text{decode}} \mid w, d] = \sum_{s \in S} R_{\text{decode}}(s, d) \cdot p(s \mid w)
\]

To account for the complete request processing pipeline, we define the overall expected throughput as the minimum of the expected prefill and decode throughputs.
Since requests cannot be served faster than the slowest stage, the effective throughput is bounded by the bottleneck.
Let $\mathbb{E}[R_{\text{prefill}} \mid w]$ denote the expected prefill throughput for workload $w$, and $\mathbb{E}[R_{\text{decode}} \mid w, d]$ the expected decode throughput for workload $w$ and average decode length $d$.
The overall expected throughput is given by:

\[
\mathbb{E}[R_{\text{overall}} \mid w, d] = \min\left( \mathbb{E}[R_{\text{prefill}} \mid w],\ \mathbb{E}[R_{\text{decode}} \mid w, d] \right)
\]

This value represents the effective throughput achievable by the system under the given workload and output length assumptions.

\begin{figure}[ht]
    \centering
    \includegraphics[width=\textwidth,height=5cm]{example-image-a}
    \caption[Sample expected throughput]
    {Sample expected throughput, calculated using the workload shown in \autoref{fig:04-sample-prefill-pdf}.}
    \label{fig:04-sample-expected-throughput}
\end{figure}

\subsection{Throughput Allocation}

Previous calculations assumed that a singular GPU island is contributing to the overall throughput of the system.
In reality, we would like to be able to construct several GPU islands and assign each island to a request type it is best suited for.
Hence, we allocate a percentage of GPU resources for a given island to a given sequence length it specialises in.
Let \(I\) be the set of all GPU islands in the system, which we partition into two separate subsets \(I_{\mathrm{prefill}}\) and \(I_{\mathrm{decode}}\).
Let \(x_i(s)\) denote the fraction of overall GPU capacity on island \(i\) devoted to handling requests of length \(s\). 
We normalise these allocations so that \(\sum_{s\in S} x_i(s) = 1 \quad\text{for every }i\in I. \)
By including \(x_i(s)\) into our previous per-length throughput models, the expected prefill throughput under workload \(w\) becomes
\[
\mathbb{E}[R_{\mathrm{prefill}}\mid w]
=\sum_{i\in I_{\mathrm{prefill}}}\sum_{s\in S}
x_i(s)\,R_{\mathrm{prefill}}^{(i)}(s)\,p(s\mid w),
\]
while the expected decode throughput for average decode length \(d\) is
\[
\mathbb{E}[R_{\mathrm{decode}}\mid w,d]
=\sum_{i\in I_{\mathrm{decode}}}\sum_{s\in S}
x_i(s)\,R_{\mathrm{decode}}^{(i)}(s,d)\,p(s\mid w).
\]

\begin{figure}[ht]
    \centering
    \includegraphics[width=\textwidth,height=5cm]{example-image-a}
    \caption[Sample throughput allocation mapping]
    {Sample throughput allocation mapping.}
    \label{fig:04-sample-throughput-allocation}
\end{figure}

% Some graphs showing the throughput calculation + Some graphs showing throughput allocation outcomes
